\chapter{系统总体设计}

\section{系统总体架构设计}
系统总体架构采用分层设计。平台划分为数据层、模型层、三维重建层、后端服务层和前端交互层。数据层负责读取原始 MRI 体数据，并完成切片提取、灰度归一化和训练集构建；模型层负责 StyleGAN2-ADA 训练、权重加载、潜变量采样和二维切片生成；三维重建层负责连续切片组织、初始点云构建、规则相机参数准备、3DGS 训练和结果导出。后端服务层基于 FastAPI 封装数据、生成、分析、任务、资产和 AI 助手接口，前端交互层则基于 React 展示仪表盘、生成工作台、医学影像查看器、报告、可信性评估和三维结果。这样的分层方式使算法计算、接口服务和页面展示之间保持清楚边界。

数据流从 IXI 数据集中的 NIfTI 体数据开始。预处理模块先把体数据转换为统一尺寸的 PNG 切片，再将切片打包为 StyleGAN2 训练数据。训练完成后，FastAPI 后端加载模型快照并生成二维 MRI 切片，React 前端通过接口获取生成结果、任务状态和图像资产。随后，系统利用潜空间插值得到连续切片序列。连续切片被转换为 PLY 点云后，再结合相机参数进入 3DGS 训练流程。训练输出包括 checkpoint、渲染图像和训练后点云，并通过后端接口提供给前端展示。

\section{数据处理模块设计}
数据处理模块的输入是 NIfTI 格式 MRI 文件，输出是用于 StyleGAN2 训练的 PNG 切片集合。该模块主要包括切片选择、灰度归一化、尺寸标准化和低信息样本过滤。

切片选择优先处理有效信息区域。系统沿体数据第三维提取轴向切片，默认保留中间 15\% 到 85\% 的深度范围，以减少头顶部、颈部和边缘空白区域的影响。灰度处理使用 2\% 和 98\% 百分位作为截断区间，再把主要灰度分布映射到 $[0,255]$，从而降低异常亮点和极端噪声的影响。尺寸标准化阶段将切片统一缩放到目标分辨率，本次 StyleGAN2 训练主要使用 $128\times128$ 图像。质量过滤阶段则根据均值、标准差和非零像素比例筛除低信息切片。

该模块的目标不是改变医学结构语义，而是让数据在格式、尺寸和灰度范围上满足生成模型训练要求。经过处理后，三维 MRI 体数据被转换为二维样本集合。该集合可用于 StyleGAN2 训练，也可用于后续结果复核。

\section{模型训练与生成模块设计}
模型训练与生成模块以 StyleGAN2-ADA 为核心。训练阶段，系统读取 \thesispath{data/processed} 目录下的数据集压缩包。系统调用 StyleGAN2 训练脚本，并传入输出目录、训练数据、GPU 数量、训练 kimg、batch size、学习率和快照保存频率等参数。训练输出包括网络快照、训练日志、训练配置文件和可视化样例。

生成阶段使用训练完成的 \texttt{network-snapshot} 权重文件。系统加载生成器的指数滑动平均模型进行推理，并根据用户输入的随机种子采样潜变量 $z$。截断参数用于调节生成结果的稳定性和多样性。一般来说，参数较小时生成结果更接近训练分布中心；参数较大时样本变化更明显，但结构也可能更不稳定。

单张生成图像不能直接表示体数据，连续切片序列由此成为二维生成和三维表达之间的中间环节。系统在两个随机种子对应的潜变量之间进行线性插值。对于序列中第 $i$ 张切片，如果设插值系数为 $t$，则潜变量可表示为：

\begin{equation}
z_i = (1-t)z_a + tz_b .
\end{equation}

通过这一设计，平台可以把离散生成样本组织为连续变化的图像序列。该序列为点云初始化和 3DGS 训练提供输入。

\section{3DGS 三维重建流程设计}
三维流程由连续切片序列构建、初始点云生成、3DGS 参数优化和结果展示四个步骤组成。这个设计来自一个直接需求：StyleGAN2 适合学习二维 MRI 切片分布，而 MRI 体数据本身具有层间顺序。换言之，系统必须在二维生成和三维重建之间建立一个可解释的中间表示。

\subsection{连续切片到空间点云的适配}
点云适配从连续切片序列开始。程序读取切片目录后，将每张图像放置到不同的 $z$ 轴位置。对于灰度值大于阈值的像素，系统把像素坐标 $(x,y)$ 和切片序号对应的 $z$ 坐标组合为空间点 $(x,y,z)$，再把像素灰度映射为 RGB 颜色。所有点最终写入 PLY 文件。

该点云反映了非背景区域在三维空间中的分布，也是 3DGS 训练和平台可视化的共同输入。PLY 文件可被常见三维工具读取，也能作为 3DGS 初始化输入。与直接生成三维体数据相比，该方式实现复杂度较低，也更容易展示二维切片进入三维空间的过程。

\subsection{三维高斯参数优化设计}
在 3DGS 重建流程中，初始点云中的空间点会被扩展为三维高斯基元。每个高斯维护中心位置 $\mu$、尺度 $s$、旋转 $r$、不透明度 $\alpha$ 和颜色 $c$ 等参数。训练目标是让高斯渲染结果接近输入切片或目标视角观测。结合 MRI 切片序列特点，模块设计时可以把像素误差、结构相似性、切片连续性和正则约束放在同一目标框架下理解：

\begin{equation}
\mathcal{L} = \lambda_1 \mathcal{L}_{pix} + \lambda_2 \mathcal{L}_{ssim} + \lambda_3 \mathcal{L}_{slice} + \lambda_4 \mathcal{L}_{reg}.
\end{equation}

其中，$\mathcal{L}_{pix}$ 约束像素差异。$\mathcal{L}_{ssim}$ 衡量结构相似性。$\mathcal{L}_{slice}$ 强化相邻切片之间的空间连续关系。$\mathcal{L}_{reg}$ 限制高斯尺度、密度和不透明度的异常变化。当前工程实验主要记录 L1 和 PSNR 指标，公式用于说明三维模块的优化思路和后续可扩展方向。

\subsection{训练接入与结果输出策略}
3DGS 训练依赖 CUDA 扩展、相机参数和显存环境。为降低平台调用复杂度，本文将三维训练封装为明确的输入输出接口。系统先完成连续切片到 PLY 点云的转换，再把点云、切片序列和相机参数文件作为训练输入。训练输出包括高斯模型、checkpoint、渲染截图、训练日志和评价指标，这些结果可直接用于平台展示和论文分析。

三维部分由此被定义为“连续切片序列组织、PLY 点云初始化、3DGS 训练优化和平台展示”的完整流程。该流程与真实项目产物对应，也使三维模块具有清楚的输入、处理和输出边界。

\section{前后端交互界面设计}
平台交互界面采用 React 实现，后端由 FastAPI 提供接口。前端按照功能场景划分为仪表盘、生成工作台、医学影像查看器、影像分析报告、可信性评估、任务中心、三维展示和教学中心等页面。其中，仪表盘用于汇总生成切片、连续序列、三维资产和风险提示；生成工作台用于选择模型权重、设置 Seed、Psi 和生成数量；医学影像查看器用于浏览 MRI 切片序列；三维展示页则用于查看体渲染结果和点云资产。

后端接口按资源和任务组织。平台通过概览接口获取系统状态，通过生成接口输出 MRI 切片，通过资产接口读取图像、模型和点云文件。任务状态由任务接口和 WebSocket 通道返回。AI 助手接口负责参数解释和教学问答。这样的接口划分使前端页面不直接依赖底层脚本，也让后端可以统一管理实验产物。

交互设计遵循流程化和低门槛原则。用户按照页面顺序完成数据查看、模型生成、影像分析、三维展示和报告整理，不需要记忆复杂命令行参数。系统通过仪表盘、任务列表、图像画廊、日志文本和三维查看器反馈运行状态，用户可以据此确认结果并整理实验材料。
